% =====================================================================
%  CHAPTER 9: 儿童少年传染病（对应大纲第九章）
%  内容来源：往届笔记第十一章，参考PPT第九章
% =====================================================================
\chapter{儿童少年传染病}
\setcounter{chapter}{9}
\chapterintro{
\textbf{其他：}儿童少年传染病流行特征、儿童少年传染病预防控制。
}

\begin{defbox}
\textbf{传染病（infectious disease）：}由病原体（细菌、病毒、立克次体、螺旋体等）引起的，能在人与人、人与动物或动物与动物间相互传染的疾病，是许多疾病的总称。
\end{defbox}

% ===== 全部内容(其他) =====
\section{儿童少年传染病流行特征}

\subsection{传染病种类}

\begin{keybox}
\textbf{1. 病毒性传染病：}流感（季节性流感）、腮腺炎、手足口病、水痘、麻疹、轮状病毒肠炎（秋季腹泻）。

\textbf{2. 细菌性传染病：}猩红热、百日咳、细菌性痢疾。
\end{keybox}

\subsection{儿童少年传染病流行状况}

\begin{keybox}
\textbf{1. 疾病分布}
\begin{enumerate}[leftmargin=*]
    \item \textbf{疾病种类分布广泛：}儿童少年中报告的传染病种类占我国法定传染病种类的80\%，其中以呼吸道疾病最多。
    \item \textbf{呼吸道传染病仍是学校传染病预防的重点：}呼吸道传染病以小学生最多，高中生和大学生次之，初中生最少；肠道传染病也是小学生最多，其余学段基本持平。
    \item \textbf{高发疾病相对集中：}季节性流感、流行性腮腺炎、手足口病、风疹等。
\end{enumerate}

\textbf{2. 人群分布}
\begin{enumerate}[leftmargin=*]
    \item \textbf{学段：}不同传播途径和高发传染病在各学段中分布不同。
    \item \textbf{性别：}学生中甲乙丙类传染病报告发病数各年龄组男生均高于女生。
    \item \textbf{年龄：}学生甲乙类传染病发病数的年龄分布呈"马鞍形"。
\end{enumerate}

\textbf{3. 地区分布：}不同省份间传染病报告发病数波动范围大。原因：（1）地区经济发展水平差异；（2）人口密度与流动情况；（3）地理气候条件（南方省份气候温暖湿润，适合蚊虫等病媒生物滋生繁殖）。

\textbf{4. 时间分布：}总体呈双峰分布，全年呈现2个发病高峰，主要集中在3-6月和10-12月。
\end{keybox}

\subsection{学校传染病流行过程特点}

\begin{keybox}
\textbf{1. 传染源的集散地}
\begin{enumerate}[leftmargin=*]
    \item \textbf{人群庞杂：}学校群体年龄结构包括儿童、青少年和成年人。
    \item \textbf{流动性强：}日常教学与社交活动导致大量跨区域、跨群体人员接触，为病原体传播创造接触网络。
    \item \textbf{健康携带者众多：}许多传染病除患者外，还有数倍甚至数十倍的健康携带者，使传染源数量、强度显著增加。
\end{enumerate}

\textbf{2. 传播特征与学生生活特点有关}
\begin{enumerate}[leftmargin=*]
    \item \textbf{群体性生活方式：}集中住宿、共用餐区和统一作息显著增加接触密度，呼吸道飞沫/接触传播风险增加。
    \item \textbf{室内微小环境中细菌聚集：}空气流通受限的室内环境中，病原体气溶胶浓度比开放空间高很多。
    \item \textbf{公共设施使用频率高：}门把手、座椅、图书馆等。
    \item \textbf{学习制度：}固定课表、集中授课等教学安排强制维持群体接触状态，日均暴露时长>6h，远超社区接触强度。
\end{enumerate}

\textbf{3. 儿童少年易感性}
\begin{enumerate}[leftmargin=*]
    \item \textbf{年龄：}年龄越小机体抵抗力越低，免疫系统作为身体抵御病原体的重要防线，婴幼儿胸腺、脾脏等免疫器官发育不完全，免疫细胞数量和活性也相对较低，使他们更易受各种传染病侵袭，如呼吸道合胞病毒感染。
    \item \textbf{特异性免疫能力：}机体特异性免疫功能尚不完善；群体内基础免疫水平不一。
    \item \textbf{卫生习惯不良。}
    \item \textbf{传染病相关知识掌握不足：}不了解传染病传播途径和预防。
    \item \textbf{特殊儿童群体免疫接种不全：}
    \begin{itemize}
        \item \textbf{城市流动儿童：}由于所在地经常变动，预防接种管理往往不规范或有所遗漏，经常发生迟种、不种、接种间隔时间不规范的情况。影响因素：人口流动性、年龄、出生地和接生场所、监护人个人素质（尤其是母亲文化程度）、预防接种服务供给利用情况。
        \item \textbf{农村留守儿童：}传染病罹患率高（健康知识缺乏、卫生习惯较差、营养摄入不足、医疗资源匮乏）；疫苗接种率低（按期接种、全程接种坚持性差，超期接种、免疫空白等现象多发）；监护人意识淡薄，交通不便，信息沟通不畅。
    \end{itemize}
\end{enumerate}
\end{keybox}

\subsection{传染病对儿童少年健康的影响}

\begin{tipbox}
\textbf{传染病对儿童少年健康的影响：}

\begin{enumerate}[leftmargin=*]
    \item \textbf{生理健康：}儿童少年罹患传染病后通常会出现发热、头痛、出疹、食欲降低等；未及时治疗会引发严重并发症。
    \item \textbf{成年期生活质量。}
    \item \textbf{疾病负担。}
\end{enumerate}
\end{tipbox}

\section{儿童少年传染病预防控制}

\begin{keybox}
\textbf{一、管理传染源}

\textbf{二、切断传播途径}

\textbf{三、保护易感人群}
\begin{enumerate}[leftmargin=*]
    \item \textbf{提高特异性免疫力：}预防接种是最经济、最有效的手段。
    \item \textbf{加强健康教育：}提高传染病防控知识和意识。
    \item \textbf{培养良好卫生习惯：}勤洗手、注意呼吸道卫生等。
    \item \textbf{改善营养和增强体质：}合理膳食，加强体育锻炼。
\end{enumerate}

\textbf{四、学校传染病防控的关键措施}
\begin{enumerate}[leftmargin=*]
    \item \textbf{晨检制度：}每日对学生进行健康检查，早期发现传染病患者。
    \item \textbf{疫情报告：}发现传染病疫情及时向疾控部门报告。
    \item \textbf{隔离治疗：}对传染病患者及时隔离，防止传播。
    \item \textbf{消毒通风：}保持教室通风，定期消毒公共设施。
    \item \textbf{健康教育：}开展传染病防治知识宣传教育。
    \item \textbf{查验接种证：}入托入学时查验预防接种证，确保免疫接种完整。
\end{enumerate}
\end{keybox}

% ----- 复习题 -----
\reviewcontent{
\section{复习题}

\begin{enumerate}[leftmargin=*, itemsep=8pt]
    \item \textbf{【名词解释】}传染病
    \item \textbf{【简答题】}简述儿童少年传染病的疾病分布特点。
    \item \textbf{【简答题】}简述学校传染病流行过程的三大特点。
    \item \textbf{【简答题】}简述儿童少年易感性的主要影响因素。
    \item \textbf{【简答题】}简述城市流动儿童和农村留守儿童免疫接种不全的原因。
    \item \textbf{【简答题】}简述学校传染病预防控制的关键措施。
\end{enumerate}

\subsection*{参考答案要点}

\begin{enumerate}[leftmargin=*, itemsep=5pt]
    \item \textbf{传染病}：由病原体（细菌、病毒、立克次体、螺旋体等）引起的，能在人与人、人与动物或动物与动物间相互传染的疾病。

    \item (1)疾病种类分布广泛，占我国法定传染病种类的80\%，以呼吸道疾病最多；(2)呼吸道传染病以小学生最多，是学校预防重点；(3)高发疾病相对集中：季节性流感、流行性腮腺炎、手足口病、风疹等；(4)时间分布呈双峰型（3-6月和10-12月）；(5)年龄分布呈"马鞍形"。

    \item (1)传染源的集散地：人群庞杂、流动性强、健康携带者众多；(2)传播特征与学生生活特点有关：群体性生活方式、室内微小环境细菌聚集、公共设施使用频率高、学习制度强制维持群体接触状态；(3)儿童少年易感性：年龄越小抵抗力越低、特异性免疫不完善、卫生习惯不良等。

    \item (1)年龄：越小抵抗力越低，免疫器官发育不完全；(2)特异性免疫能力不完善，群体内免疫水平不一；(3)卫生习惯不良；(4)传染病相关知识掌握不足；(5)特殊群体免疫接种不全（城市流动儿童和农村留守儿童）。

    \item 城市流动儿童：(1)人口流动性（居住时间、异地来往频率）；(2)年龄越小完成率越高；(3)出生地和接生场所；(4)监护人素质（尤其母亲文化程度）；(5)预防接种服务供给不均衡。农村留守儿童：(1)监护人意识淡薄，交通不便，信息不畅；(2)健康知识缺乏，营养不足，医疗资源匮乏；(3)超期接种、免疫空白现象多发。

    \item (1)晨检制度，早期发现；(2)疫情报告，及时向疾控部门报告；(3)隔离治疗，防止传播；(4)消毒通风，保持空气流通；(5)健康教育，提高防控意识；(6)查验接种证，确保免疫完整。预防接种是最经济、最有效的手段。
\end{enumerate}
}

\newpage
